\section{Conclusiones}
\label{sec:conclusiones}

\paragraph*{Diseño} El circuito obtiene las cuatro señales de la consigna en el orden pedido y cada una a partir de la anterior, con tres TL084 y $\pm\SI{12}{\volt}$. El TL084 se eligió por su entrada JFET, que hace despreciable el error por corriente de polarización en el integrador, por su \emph{slew rate}, que el rectificador necesita para conmutar dentro de impulsos de pocos microsegundos, y porque tiene cuatro amplificadores por integrado. La impedancia de entrada es de \SI{15}{\kilo\ohm}, fijada por un resistor y un seguidor, y la de salida es de \valu{2}{\simZoutF}{\ohm} a la frecuencia de trabajo; como es inductiva ($L_{eq} \approx \valu{0}{\teoLeq}{\micro\henry}$), se mide inyectando corriente y no con el método de la carga.

\paragraph*{Niveles de continua} Los niveles salen de un zener con presets y seguidores, y no de divisores sobre la fuente, y ninguna etapa amplifica la continua con ganancia mayor que 2: cada amplitud se fija en la etapa que genera la señal. El punto crítico fue el integrador. Con una resistencia en paralelo hay que elegir entre rampas rectas y un nivel estable: con $R_F/R_I = 100$, \SI{50}{\milli\volt} de error en la entrada mueven la triangular \valu{1}{\simIntDcCienCorr}{\volt}. El capacitor en serie con $R_I$ resuelve ese compromiso: la triangular queda centrada en la referencia de \SI{1,25}{\volt}, como pide la consigna, con una alinealidad del \valu{2}{\simIntAcNoLin}{\percent}.

\paragraph*{Estabilidad y velocidad} La señal 3 se obtuvo con un derivador y no con un comparador, porque un comparador contra el nivel medio la habría atrasado un cuarto de período. Sin compensar, ese derivador oscila a unos \SI{130}{\kilo\hertz}; con $R_S = \SI{1}{\kilo\ohm}$ y $C_F = \SI{56}{\pico\farad}$ el sobrepico baja al \valu{1}{\simBarrDSobrepico}{\percent}, al costo de flancos de unos \SI{4}{\micro\second}. El rectificador de precisión funciona con 1N4004, aunque con 1N4148 el pulso negativo de conmutación sería la mitad.

\paragraph*{Implementación} En la placa, la señal 1 coincidió con la simulación hecha con la entrada medida, y la triangular tuvo la forma, la fase y la linealidad del diseño.

\paragraph*{Lo que salió mal} Los niveles de las señales 1 y 2 quedaron corridos porque se midió sin terminar de ajustar los presets, y la triangular salió un \valu{0}{100*(\medSDosVertVpp/\simMedSdosVpp-1)}{\percent} más grande porque el capacitor de integración se apartaba de su valor nominal. La señal 4 tuvo un solo impulso dominante por período y ráfagas de oscilación. Las capturas no alcanzaron para ubicar la causa: hay un solo canal por captura, no se registró la señal 3 y la base de tiempo no resolvía la oscilación. La simulación muestra que el diseño nominal no produce esos síntomas y que los reproducen un derivador poco amortiguado y un desbalance en la red sumadora del rectificador; las mediciones para confirmarlo están al final de la Sección~\ref{sec:fallas}.

\paragraph*{Aprendizajes} La simulación con el macromodelo del fabricante fue indispensable en las dos etapas que dependen de la dinámica del operacional, el derivador y el rectificador: con un operacional ideal ninguna de las dos muestra problemas. También lo fue para interpretar las mediciones, porque comparando las capturas con el circuito simulado, con y sin defectos, se pudo descartar un error de diseño y acotar dónde buscar. De la medición queda una lección práctica: ajustar los presets antes de capturar, medir con dos canales para tener referencia de fase y elegir la base de tiempo según el detalle que se quiere ver, no según el período de la señal.
